% Generated draft: requires mathematical and source review.
\documentclass[11pt]{article}
\usepackage[margin=25mm]{geometry}
\usepackage{fontspec}
\setmainfont{Latin Modern Roman}
\usepackage{amsmath,amssymb,mathtools,unicode-math}
\setmathfont{Latin Modern Math}
\usepackage{enumitem,needspace,xurl,hyperref}
\setlength{\parindent}{0pt}
\setlength{\parskip}{5pt}
\setlength{\emergencystretch}{3em}
\newcommand{\N}{\mathbb{N}}
\newcommand{\Z}{\mathbb{Z}}
\newcommand{\Q}{\mathbb{Q}}
\newcommand{\R}{\mathbb{R}}
\newcommand{\C}{\mathbb{C}}
\newcommand{\F}{\mathbb{F}}
\newcommand{\A}{\mathbb{A}}
\newcommand{\PP}{\mathbb{P}}
\newcommand{\E}{\mathbb{E}}
\newcommand{\eps}{\varepsilon}
\newcommand{\dd}{\,\mathrm{d}}
\begin{document}
\section*{UnsolvedMath ID: 801143}
\textbf{Source problem code:} catalogue entry
\subsection{Definitions and mathematical statement}
In the usual complex setting, a type \(\mathrm{II}_1\) factor is a von Neumann algebra \(P\) with center \(Z(P)=\C 1\), no nonzero minimal projections, and a faithful normal tracial state \(\tau_P\) satisfying \(\tau_P(1)=1\). Let \(\operatorname{Th}_{\mathrm{tr}}(P,\tau_P)\) denote its continuous first-order theory in the standard language of tracial von Neumann algebras, using operator-norm bounded sorts and the metric \(d_2(x,y)=\tau_P((x-y)^*(x-y))^{1/2}\). Let \(\operatorname{Th}_{C^*}(P)\) denote the continuous first-order theory of its underlying unital \(C^*\)-algebra in the standard norm language, without a distinguished trace. Equality of theories means equality of the values of all sentences in the respective language.

Under the interpretation that the first occurrence of ``elementarily equivalent'' refers to the tracial von Neumann algebra language, and that ``but not as as \(C^*\)-algebras'' means failure of elementary equivalence in the \(C^*\)-algebra language, the question is whether
\[
\exists\, N,M\text{ of type }\mathrm{II}_1:\qquad
\operatorname{Th}_{\mathrm{tr}}(N,\tau_N)=\operatorname{Th}_{\mathrm{tr}}(M,\tau_M)
\quad\text{and}\quad
\operatorname{Th}_{C^*}(N)\ne\operatorname{Th}_{C^*}(M).
\]
This is a conditional reconstruction of the supplied wording. The source does not explicitly specify the first language, and its phrase ``not as as'' is defective. It also refers to a definition of \(2\)-quantifier theory at \texttt{[quantifiertheory]}, but supplies neither that definition nor the referenced bibliographic entry. Whether a restriction to \(2\)-quantifier theory was intended remains unresolved; the displayed reconstruction uses full elementary equivalence.
\subsection{Short English statement}
Under the stated interpretation, the question asks whether two type \(\mathrm{II}_1\) factors can have the same first-order theory as tracial von Neumann algebras but different first-order theories as unital \(C^*\)-algebras. The intended languages, the defective wording, and the role of the reference to \(2\)-quantifier theory require confirmation.
\subsection{Sources}
\begin{itemize}
\item Local UnsolvedMath registry: permanent ID \(801143\), source problem code \texttt{catalogue entry}, title ``Ultraproduct \(\mathrm{II}_1\) Factors --- Problem 45.'' The supplied record attributes the problem to Elayavalli, with proposed year 2025, and labels it open. This status is supplied registry information, not an independently verified mathematical assessment.
\item Supplied source attribution: Elayavalli, \emph{50 Open Problems: Ultraproduct \(\mathrm{II}_1\) Factors} (2025), Problem 45, \url{https://arxiv.org/abs/2511.20377}. The record reports access on 2026-10-16 and extraction by \texttt{arxiv\_latex}; these provenance details are supplied and were not independently verified here.
\item The statement cites \texttt{[quantifiertheory]} for the definition of \(2\)-quantifier theory. No corresponding bibliographic entry, definition, or link is supplied.
\end{itemize}
\end{document}
